\begin{pastexam}[label={px:ust-115-ee-ds-36}]{臺聯大電機類 115 學年度 第 36 題}
（單選；2 分）

Which of the following C functions/macros are valid codes to swap two int variables?
\begin{minted}[fontsize=\small]{c}
M1. #define SWAP(type,a,b) {type temp=a;a=b;b=temp;}
    //SWAP(int,a,b)
M2. #define SWAP(a,b) (a+=b,b=a-b,a-=b)
    //SWAP(a,b);
F3. void swap(int* a, int* b) {
        int temp;
        temp = *a;
        *a = *b;
        *b = temp;
    }
    //int x,y; swap(&x,&y);
F4. void swap(int* a, int* b) {
        *a = *b;
        *b = *a;
    }
    //int x,y; swap(&x,&y);
F5. void swap(int a, int b) {
        int temp;
        temp = a;
        a = b;
        b = temp;
        Return;
    }
    //int x,y; swap(x,y);
\end{minted}

(A)~M1 and M2 only \qquad (B)~F3 only \qquad (C)~M2 and F3 \qquad (D)~M1, M2, and F3 \qquad (E)~All
\end{pastexam}

\begin{pastsoln}
答案是 (D)。

逐一檢查（設 \texttt{a}、\texttt{b} 是兩個不同的 \texttt{int} 變數）：
\begin{itemize}
\item \textbf{M1}：\texttt{SWAP(int,a,b)} 展開成 \texttt{\{int temp=a;a=b;b=temp;\}}，是一個區塊，
用暫存變數交換，正確。 （小缺點：如果要交換的變數本身叫 \texttt{temp} 會出錯；在 \texttt{if-else}
裡寫成 \texttt{SWAP(...);} 時多出的分號會讓 \texttt{else} 對不上，所以一般寫成
\texttt{do \{ ... \} while (0)}。 照題目給的用法沒有問題。）
\item \textbf{M2}：設 \texttt{a} $= x$、\texttt{b} $= y$。 \texttt{a+=b} 後 \texttt{a} $= x + y$；
\texttt{b=a-b} 後 \texttt{b} $= x$；\texttt{a-=b} 後 \texttt{a} $= y$。 逗號運算子保證由左到右
依序執行，正確。 （限制：$x + y$ 溢位時是未定義行為；\texttt{SWAP(x,x)} 會把 \texttt{x} 變成 $0$。
題目問一般的兩個變數，視為正確。）
\item \textbf{F3}：傳入兩個位址，透過指標用暫存變數交換，是標準寫法，正確。
\item \textbf{F4}：\texttt{*a = *b} 先把 \texttt{*a} 蓋掉，原來的值就沒了；\texttt{*b = *a} 只是把
\texttt{*b} 指定成自己。 結果兩個都變成原來的 \texttt{y}。 錯。
\item \textbf{F5}：參數是傳值，函式裡交換的是 \texttt{a}、\texttt{b} 這兩個複本，呼叫者的
\texttt{x}、\texttt{y} 不變。 而且 C 區分大小寫，\texttt{Return;} 不是 \texttt{return}，連編譯都
不會過。 錯。
\end{itemize}
正確的是 M1、M2、F3，選 (D)。 (A) 漏了 F3，(B) 漏了 M1、M2，(C) 漏了 M1，(E) 把 F4、F5 也算進去。
\end{pastsoln}
